\documentclass[11pt,a4paper]{article}

\usepackage[margin=1in]{geometry}
\usepackage{amsmath,amssymb}
\usepackage{booktabs}
\usepackage{xcolor}
\usepackage{listings}
\usepackage{hyperref}
\usepackage{titlesec}
\usepackage{microtype}
\usepackage{parskip}
clear
clear
clear
exit



% ---------------------------------------------------------------------------
% Colour palette
% ---------------------------------------------------------------------------
\definecolor{codebg}{RGB}{245,246,248}
\definecolor{kw}{RGB}{0,80,160}
\definecolor{cm}{RGB}{60,128,60}
\definecolor{st}{RGB}{160,40,40}
\definecolor{nu}{RGB}{140,60,180}
\definecolor{pp}{RGB}{100,50,0}
\definecolor{rulecol}{RGB}{180,190,200}

% ---------------------------------------------------------------------------
% C++20 listing style
% ---------------------------------------------------------------------------
\lstdefinestyle{cpp20}{
  language=C++,
  backgroundcolor=\color{codebg},
  basicstyle=\ttfamily\small,
  keywordstyle=\color{kw}\bfseries,
  commentstyle=\color{cm}\itshape,
  stringstyle=\color{st},
  numberstyle=\color{nu}\tiny,
  emphstyle=\color{pp}\bfseries,
  emph={constexpr,noexcept,nodiscard,string_view,int32_t,
        size_t,uint32_t,ParticleID,DecayEvent,KineticEventRecord,
        EventKind,EventPayload,NuclearState,DaughterState,
        GenericTransportEvent},
  breaklines=true,
  frame=single,
  rulecolor=\color{rulecol},
  framesep=4pt,
  numbers=left,
  numbersep=8pt,
  tabsize=4,
  showstringspaces=false,
  xleftmargin=14pt,
  captionpos=b
}

% C style
\lstdefinestyle{clang}{
  language=C,
  backgroundcolor=\color{codebg},
  basicstyle=\ttfamily\small,
  keywordstyle=\color{kw}\bfseries,
  commentstyle=\color{cm}\itshape,
  stringstyle=\color{st},
  numberstyle=\color{nu}\tiny,
  breaklines=true,
  frame=single,
  rulecolor=\color{rulecol},
  framesep=4pt,
  numbers=left,
  numbersep=8pt,
  tabsize=4,
  showstringspaces=false,
  xleftmargin=14pt,
  captionpos=b
}

\lstset{style=cpp20}

% ---------------------------------------------------------------------------
% Section formatting
% ---------------------------------------------------------------------------
\titleformat{\section}{\large\bfseries}{}{0em}{}[\titlerule]
\titleformat{\subsection}{\normalsize\bfseries}{}{0em}{}

% ---------------------------------------------------------------------------
\begin{document}

\begin{center}
  {\LARGE\bfseries Particle Identifier and Decay Event Subsystem}\\[4pt]
  {\large VSEPR-SIM --- Code Reference Draft}\\[2pt]
  {\small\texttt{include/physics/} \quad \texttt{include/kinetics/} \quad
         \texttt{legacy/c\_api/} \quad \texttt{tests/}}\\[6pt]
  {\small v4.0.0 \quad June 2025}
\end{center}

\vspace{6pt}
\hrule
\vspace{10pt}

\tableofcontents
\newpage

% ============================================================================
\section{Overview}
% ============================================================================

The particle identifier subsystem defines a signed-integer identity
namespace that spans ordinary matter species, emitted decay particles,
radiation bookkeeping tokens, and defect physics sentinels.  Every entity
the engine emits, tracks, or accounts for is tagged with a value of type
\texttt{ParticleID}.

\[
ID_p \in \mathbb{Z}, \quad
\begin{cases}
ID_p \ge 0 & \text{ordinary matter (atoms, isotopes, molecules, candidates)} \\
ID_p  <  0 & \text{emitted particles, energy stand-ins, defect tokens}
\end{cases}
\]

The entire scheme is constexpr-queryable at compile time.  No runtime
tables, no lookup files, no hidden state.

\bigskip
\begin{center}
\begin{tabular}{lll}
\toprule
\textbf{Header} & \textbf{Namespace} & \textbf{Contents} \\
\midrule
\texttt{include/physics/particle\_id.hpp}           & \texttt{vsepr::physics} & \texttt{ParticleID} enum + 5 constexpr queries \\
\texttt{include/physics/decay\_event.hpp}            & \texttt{vsepr::physics} & \texttt{DecayEvent} struct \\
\texttt{include/physics/decay\_event\_builder.hpp}   & \texttt{vsepr::physics} & Builder factories + daughter arithmetic \\
\texttt{include/kinetics/kinetic\_event\_kind.hpp}   & \texttt{vsepr::kinetics} & \texttt{EventKind} enum \\
\texttt{include/kinetics/kinetic\_event\_record.hpp} & \texttt{vsepr::kinetics} & \texttt{KineticEventRecord} + \texttt{EventPayload} \\
\texttt{legacy/c\_api/particle\_ids.h/.c}            & --- (plain C) & C enum + 4 query functions \\
\texttt{legacy/c\_api/decay\_event.h/.c}             & --- (plain C) & C struct + init/add \\
\bottomrule
\end{tabular}
\end{center}

% ============================================================================
\section{\texttt{particle\_id.hpp} --- Identity Layer}
\label{sec:particle-id}
% ============================================================================

\subsection{The \texttt{ParticleID} Enum}

\texttt{ParticleID} is a scoped enum with \texttt{std::int32\_t} backing.
Negative values are reserved.  Positive values are assigned at runtime by
the species registry.

\begin{lstlisting}[style=cpp20, caption={\texttt{include/physics/particle\_id.hpp} --- enum definition}]
namespace vsepr::physics {

enum class ParticleID : std::int32_t {
    nuclear_placeholder    =  0,

    // Emitted decay particles
    alpha_particle         = -1,
    beta_minus             = -2,
    gamma_photon           = -3,
    energy_packet_generic  = -4,
    beta_plus              = -5,
    electron_capture_token = -6,
    neutron_emission       = -7,
    proton_emission        = -8,
    neutrino_loss_token    = -9,

    // Defect tokens
    vacancy_token          = -10,
    interstitial_token     = -11,
    defect_cluster_token   = -12,
    helium_bubble_seed     = -13,
    lattice_damage_packet  = -14
};

} // namespace vsepr::physics
\end{lstlisting}

The enum partitions into three sub-ranges that the constexpr queries
exploit directly:

\[
\underbrace{[-14,\,-10]}_{\text{defect tokens}}
\quad
\underbrace{[-9,\,-1]}_{\text{emitted particles}}
\quad
\underbrace{\{0\}}_{\text{placeholder}}
\quad
\underbrace{[1,\,+\infty)}_{\text{matter species (runtime)}}
\]

\subsection{Constexpr Query Functions}

All five queries are \texttt{[[nodiscard]] constexpr noexcept}. They cost
zero runtime cycles when the argument is a compile-time constant.

\begin{lstlisting}[style=cpp20, caption={Constexpr queries --- \texttt{particle\_id.hpp}}]
// Reserved: 0 through -99 (the full engine sentinel range)
[[nodiscard]] constexpr bool is_reserved_particle_id(ParticleID id) noexcept {
    const auto v = static_cast<std::int32_t>(id);
    return (v <= 0 && v >= -99);
}

// Transportable: physically propagates through matter
[[nodiscard]] constexpr bool is_transportable(ParticleID id) noexcept {
    switch (id) {
        case ParticleID::alpha_particle:
        case ParticleID::beta_minus:
        case ParticleID::gamma_photon:
        case ParticleID::beta_plus:
        case ParticleID::neutron_emission:
        case ParticleID::proton_emission:
            return true;
        default:
            return false;
    }
}

// Bookkeeping only: token exists for accounting, never transported
[[nodiscard]] constexpr bool is_bookkeeping_only(ParticleID id) noexcept {
    switch (id) {
        case ParticleID::energy_packet_generic:
        case ParticleID::electron_capture_token:
        case ParticleID::neutrino_loss_token:
            return true;
        default:
            return false;
    }
}

// Defect token: range check [-14, -10]
[[nodiscard]] constexpr bool is_defect_token(ParticleID id) noexcept {
    const auto v = static_cast<std::int32_t>(id);
    return (v <= -10 && v >= -14);
}

// String label for logging and output
[[nodiscard]] constexpr std::string_view to_string(ParticleID id) noexcept {
    switch (id) {
        case ParticleID::nuclear_placeholder:    return "nuclear_placeholder";
        case ParticleID::alpha_particle:         return "alpha_particle";
        case ParticleID::beta_minus:             return "beta_minus";
        case ParticleID::gamma_photon:           return "gamma_photon";
        case ParticleID::energy_packet_generic:  return "energy_packet_generic";
        case ParticleID::beta_plus:              return "beta_plus";
        case ParticleID::electron_capture_token: return "electron_capture_token";
        case ParticleID::neutron_emission:       return "neutron_emission";
        case ParticleID::proton_emission:        return "proton_emission";
        case ParticleID::neutrino_loss_token:    return "neutrino_loss_token";
        case ParticleID::vacancy_token:          return "vacancy_token";
        case ParticleID::interstitial_token:     return "interstitial_token";
        case ParticleID::defect_cluster_token:   return "defect_cluster_token";
        case ParticleID::helium_bubble_seed:     return "helium_bubble_seed";
        case ParticleID::lattice_damage_packet:  return "lattice_damage_packet";
        default:                                 return "unknown_particle_id";
    }
}
\end{lstlisting}

\subsection{Query Classification Table}

\begin{center}
\begin{tabular}{lcccc}
\toprule
\textbf{ParticleID} & \textbf{reserved} & \textbf{transportable} & \textbf{bookkeeping} & \textbf{defect} \\
\midrule
\texttt{nuclear\_placeholder}    & \checkmark & & & \\
\texttt{alpha\_particle}         & \checkmark & \checkmark & & \\
\texttt{beta\_minus}             & \checkmark & \checkmark & & \\
\texttt{gamma\_photon}           & \checkmark & \checkmark & & \\
\texttt{energy\_packet\_generic} & \checkmark & & \checkmark & \\
\texttt{beta\_plus}              & \checkmark & \checkmark & & \\
\texttt{electron\_capture\_token}& \checkmark & & \checkmark & \\
\texttt{neutron\_emission}       & \checkmark & \checkmark & & \\
\texttt{proton\_emission}        & \checkmark & \checkmark & & \\
\texttt{neutrino\_loss\_token}   & \checkmark & & \checkmark & \\
\texttt{vacancy\_token}          & \checkmark & & & \checkmark \\
\texttt{interstitial\_token}     & \checkmark & & & \checkmark \\
\texttt{defect\_cluster\_token}  & \checkmark & & & \checkmark \\
\texttt{helium\_bubble\_seed}    & \checkmark & & & \checkmark \\
\texttt{lattice\_damage\_packet} & \checkmark & & & \checkmark \\
\bottomrule
\end{tabular}
\end{center}

% ============================================================================
\section{\texttt{decay\_event.hpp} --- Event Record}
\label{sec:decay-event}
% ============================================================================

\texttt{DecayEvent} is a plain aggregate that captures the full
information of a single nuclear transition.  Default member
initialisation means a value-constructed \texttt{DecayEvent\{\}} is
always zero-safe.

\begin{lstlisting}[style=cpp20, caption={\texttt{include/physics/decay\_event.hpp}}]
namespace vsepr::physics {

struct DecayEvent {
    static constexpr std::size_t max_emitted_particles = 8;

    std::int32_t parent_species_id  {0};
    std::int32_t daughter_species_id {0};

    std::array<ParticleID, max_emitted_particles> emitted_particle_ids {};
    std::size_t emitted_count {0};

    double released_energy_eV          {0.0};
    double deposited_energy_fraction   {0.0};
    double transported_energy_fraction {0.0};
    double local_damage_score          {0.0};

    double event_time_s {0.0};
    double confidence   {0.0};

    [[nodiscard]] bool add_emitted(ParticleID id) noexcept {
        if (emitted_count >= emitted_particle_ids.size()) {
            return false;
        }
        emitted_particle_ids[emitted_count++] = id;
        return true;
    }
};

} // namespace vsepr::physics
\end{lstlisting}

\subsection{Field Semantics}

\begin{center}
\begin{tabular}{lll}
\toprule
\textbf{Field} & \textbf{Type} & \textbf{Meaning} \\
\midrule
\texttt{parent\_species\_id}            & \texttt{int32\_t} & Engine-assigned ID of parent nuclide \\
\texttt{daughter\_species\_id}          & \texttt{int32\_t} & Engine-assigned ID of daughter nuclide \\
\texttt{emitted\_particle\_ids}         & \texttt{array<ParticleID,8>} & Reserved IDs of emitted particles \\
\texttt{emitted\_count}                 & \texttt{size\_t}  & How many slots are populated \\
\texttt{released\_energy\_eV}           & \texttt{double}   & Total Q-value in electron-volts \\
\texttt{deposited\_energy\_fraction}    & \texttt{double}   & Fraction absorbed locally $\in [0,1]$ \\
\texttt{transported\_energy\_fraction}  & \texttt{double}   & Fraction carried away $\in [0,1]$ \\
\texttt{local\_damage\_score}           & \texttt{double}   & Lattice displacement proxy $\in [0,1]$ \\
\texttt{event\_time\_s}                 & \texttt{double}   & Engine time of event in seconds \\
\texttt{confidence}                     & \texttt{double}   & Source reliability $\in [0,1]$ \\
\bottomrule
\end{tabular}
\end{center}

The invariant \texttt{deposited\_energy\_fraction + transported\_energy\_fraction $\approx$ 1}
is enforced by the builder factories but not by the struct itself, which
remains a plain data holder.

% ============================================================================
\section{\texttt{decay\_event\_builder.hpp} --- Factory Layer}
\label{sec:builder}
% ============================================================================

\subsection{Nuclear State and Daughter Arithmetic}

\texttt{NuclearState} is a minimal struct carrying only the fields
needed for transition arithmetic.  It does \emph{not} replace
\texttt{ElementDescriptor} (the full identity layer).

\begin{lstlisting}[style=cpp20, caption={Transition structs and constexpr daughter arithmetic}]
namespace vsepr::physics {

struct NuclearState {
    std::int32_t species_id {0};
    int Z {0};
    int A {0};
};

struct DaughterState {
    int Z {0};
    int A {0};
};

[[nodiscard]] constexpr DaughterState
alpha_daughter(const NuclearState& s) noexcept {
    return DaughterState{ s.Z - 2, s.A - 4 };
}

[[nodiscard]] constexpr DaughterState
beta_minus_daughter(const NuclearState& s) noexcept {
    return DaughterState{ s.Z + 1, s.A };
}

[[nodiscard]] constexpr DaughterState
beta_plus_daughter(const NuclearState& s) noexcept {
    return DaughterState{ s.Z - 1, s.A };
}

} // namespace vsepr::physics
\end{lstlisting}

The arithmetic corresponds directly to the standard nuclear transition
equations:
\[
\alpha:\quad S(Z,A) \;\to\; S(Z-2,\,A-4) + P_{-1}
\]
\[
\beta^-:\quad S(Z,A) \;\to\; S(Z+1,\,A) + P_{-2} + P_{-9}
\]
\[
\beta^+:\quad S(Z,A) \;\to\; S(Z-1,\,A) + P_{-5} + P_{-9}
\]

\subsection{Builder Factories}

Each factory populates a \texttt{DecayEvent} with mode-specific energy
partitioning defaults and the correct emitted-particle list.

\begin{lstlisting}[style=cpp20, caption={Alpha decay factory}]
// Alpha: heavy, short-range, high local damage
// Deposited 0.85 / Transported 0.15 / Damage 0.90
[[nodiscard]] inline DecayEvent make_alpha_decay(
    std::int32_t parent_species_id,
    std::int32_t daughter_species_id,
    double released_energy_eV,
    double event_time_s,
    double confidence
) {
    DecayEvent ev {};
    ev.parent_species_id   = parent_species_id;
    ev.daughter_species_id = daughter_species_id;
    ev.released_energy_eV  = released_energy_eV;
    ev.event_time_s        = event_time_s;
    ev.confidence          = confidence;
    ev.deposited_energy_fraction   = 0.85;
    ev.transported_energy_fraction = 0.15;
    ev.local_damage_score          = 0.90;
    (void)ev.add_emitted(ParticleID::alpha_particle);
    return ev;
}
\end{lstlisting}

\begin{lstlisting}[style=cpp20, caption={Beta-minus decay factory}]
// Beta-minus: light, medium range, moderate damage
// Deposited 0.35 / Transported 0.65 / Damage 0.30
[[nodiscard]] inline DecayEvent make_beta_minus_decay(
    std::int32_t parent_species_id,
    std::int32_t daughter_species_id,
    double released_energy_eV,
    double event_time_s,
    double confidence
) {
    DecayEvent ev {};
    ev.parent_species_id   = parent_species_id;
    ev.daughter_species_id = daughter_species_id;
    ev.released_energy_eV  = released_energy_eV;
    ev.event_time_s        = event_time_s;
    ev.confidence          = confidence;
    ev.deposited_energy_fraction   = 0.35;
    ev.transported_energy_fraction = 0.65;
    ev.local_damage_score          = 0.30;
    (void)ev.add_emitted(ParticleID::beta_minus);
    (void)ev.add_emitted(ParticleID::neutrino_loss_token);
    return ev;
}
\end{lstlisting}

\begin{lstlisting}[style=cpp20, caption={Gamma decay factory}]
// Gamma: penetrating, low local damage, high transport
// Deposited 0.10 / Transported 0.90 / Damage 0.05
[[nodiscard]] inline DecayEvent make_gamma_decay(
    std::int32_t parent_species_id,
    std::int32_t daughter_species_id,
    double released_energy_eV,
    double event_time_s,
    double confidence
) {
    DecayEvent ev {};
    ev.parent_species_id   = parent_species_id;
    ev.daughter_species_id = daughter_species_id;
    ev.released_energy_eV  = released_energy_eV;
    ev.event_time_s        = event_time_s;
    ev.confidence          = confidence;
    ev.deposited_energy_fraction   = 0.10;
    ev.transported_energy_fraction = 0.90;
    ev.local_damage_score          = 0.05;
    (void)ev.add_emitted(ParticleID::gamma_photon);
    return ev;
}
\end{lstlisting}

\subsection{Energy Partitioning Defaults}

\begin{center}
\begin{tabular}{lccc}
\toprule
\textbf{Factory} & \textbf{Deposited} & \textbf{Transported} & \textbf{Damage score} \\
\midrule
\texttt{make\_alpha\_decay}     & 0.85 & 0.15 & 0.90 \\
\texttt{make\_beta\_minus\_decay} & 0.35 & 0.65 & 0.30 \\
\texttt{make\_gamma\_decay}     & 0.10 & 0.90 & 0.05 \\
\bottomrule
\end{tabular}
\end{center}

% ============================================================================
\section{\texttt{kinetic\_event\_kind.hpp} --- Classification}
\label{sec:eventkind}
% ============================================================================

\texttt{EventKind} is the coarse universal classifier that sits at the top
of the kinetics layer.  Formation-specific \texttt{EventType} values from
\texttt{formation\_kinetics.hpp} (28-type taxonomy) map into the 11 kinds
here at the scheduling layer.

\begin{lstlisting}[style=cpp20, caption={\texttt{include/kinetics/kinetic\_event\_kind.hpp}}]
namespace vsepr::kinetics {

enum class EventKind : std::int32_t {
    unknown        = 0,
    decay,
    transport,
    collision,
    adsorption,
    desorption,
    coalescence,
    nucleation,
    deposition,
    restructuring,
    phase_change
};

} // namespace vsepr::kinetics
\end{lstlisting}

% ============================================================================
\section{\texttt{kinetic\_event\_record.hpp} --- Kinetics Wrapper}
\label{sec:record}
% ============================================================================

\texttt{KineticEventRecord} wraps any physics-specific event in a single
schedulable unit.  The payload is a \texttt{std::variant} so the kinetic
scheduler never needs to cast or union-punch.

\begin{lstlisting}[style=cpp20, caption={\texttt{include/kinetics/kinetic\_event\_record.hpp}}]
namespace vsepr::kinetics {

struct GenericTransportEvent {
    int source_id {0};
    int target_id {0};
    double time_s {0.0};
    double confidence {0.0};
};

using EventPayload = std::variant<
    vsepr::physics::DecayEvent,
    GenericTransportEvent
>;

struct KineticEventRecord {
    EventKind kind {EventKind::unknown};
    double intensity {0.0};
    double barrier_score {0.0};
    double environment_multiplier {1.0};
    EventPayload payload {};
};

} // namespace vsepr::kinetics
\end{lstlisting}

\subsection{Dispatch Pattern}

The canonical dispatch pattern over \texttt{EventPayload} uses
\texttt{std::get\_if}, which is zero-overhead on most ABIs:

\begin{lstlisting}[style=cpp20, caption={Variant dispatch at the scheduling layer}]
void dispatch(const KineticEventRecord& rec) {
    if (auto* d = std::get_if<vsepr::physics::DecayEvent>(&rec.payload)) {
        // nuclear decay path
        process_decay(*d);
    } else if (auto* t = std::get_if<GenericTransportEvent>(&rec.payload)) {
        // transport path
        process_transport(*t);
    }
}
\end{lstlisting}

% ============================================================================
\section{Plain C Legacy API}
\label{sec:c-api}
% ============================================================================

The plain C bridge in \texttt{legacy/c\_api/} mirrors the C++ enum with
SCREAMING\_SNAKE\_CASE constants and provides four query functions via
\texttt{extern "C"} linkage.

\begin{lstlisting}[style=clang, caption={\texttt{legacy/c\_api/particle\_ids.h}}]
typedef enum ParticleID {
    PARTICLE_NUCLEAR_PLACEHOLDER    =  0,

    PARTICLE_ALPHA                  = -1,
    PARTICLE_BETA_MINUS             = -2,
    PARTICLE_GAMMA                  = -3,
    PARTICLE_ENERGY_PACKET          = -4,
    PARTICLE_BETA_PLUS              = -5,
    PARTICLE_ELECTRON_CAPTURE       = -6,
    PARTICLE_NEUTRON_EMISSION       = -7,
    PARTICLE_PROTON_EMISSION        = -8,
    PARTICLE_NEUTRINO_LOSS          = -9,

    PARTICLE_VACANCY                = -10,
    PARTICLE_INTERSTITIAL           = -11,
    PARTICLE_DEFECT_CLUSTER         = -12,
    PARTICLE_HELIUM_BUBBLE_SEED     = -13,
    PARTICLE_LATTICE_DAMAGE_PACKET  = -14
} ParticleID;

const char* particle_id_name(int id);
int particle_id_is_reserved(int id);
int particle_id_is_transportable(int id);
int particle_id_is_bookkeeping_only(int id);
\end{lstlisting}

\begin{lstlisting}[style=clang, caption={\texttt{legacy/c\_api/decay\_event.h}}]
#define MAX_EMITTED_PARTICLES 8

typedef struct DecayEvent {
    int parent_species_id;
    int daughter_species_id;

    int emitted_particle_ids[MAX_EMITTED_PARTICLES];
    size_t emitted_count;

    double released_energy_eV;
    double deposited_energy_fraction;
    double transported_energy_fraction;
    double local_damage_score;

    double event_time_s;
    double confidence;
} DecayEvent;

void decay_event_init(DecayEvent* ev);
int decay_event_add_emitted(DecayEvent* ev, int particle_id);
\end{lstlisting}

% ============================================================================
\section{Test Suite --- \texttt{test\_particle\_id\_decay.cpp}}
\label{sec:tests}
% ============================================================================

The test executable runs 76 assertions across 8 groups and exits
non-zero on any failure.

\begin{lstlisting}[style=cpp20, caption={Test harness macro and structure}]
#define CHECK(cond, msg) do { \
    ++tests_run; \
    if (cond) { ++tests_passed; std::printf("  [PASS] %s\n", msg); } \
    else      { std::printf("  [FAIL] %s  (line %d)\n", msg, __LINE__); } \
} while(0)
\end{lstlisting}

\begin{lstlisting}[style=cpp20, caption={Alpha decay test --- Pu-239 to U-235}]
auto ev = make_alpha_decay(
    /*parent_species_id=*/  942390,
    /*daughter_species_id=*/922350,
    /*released_energy_eV=*/ 5.24e6,
    /*event_time_s=*/        1.0,
    /*confidence=*/          0.98
);

CHECK(ev.parent_species_id   == 942390, "parent ID");
CHECK(ev.daughter_species_id == 922350, "daughter ID");
CHECK(ev.emitted_count == 1,            "one emitted particle");
CHECK(ev.emitted_particle_ids[0] == ParticleID::alpha_particle, "emitted alpha");
CHECK(ev.released_energy_eV > 5.0e6,   "energy > 5 MeV");
CHECK(ev.deposited_energy_fraction > 0.8,   "high deposition");
CHECK(ev.local_damage_score > 0.8,          "high damage");
\end{lstlisting}

\begin{lstlisting}[style=cpp20, caption={Overflow guard --- capacity capped at 8}]
DecayEvent ev {};
for (int i = 0; i < 8; ++i) {
    CHECK(ev.add_emitted(ParticleID::gamma_photon), "add succeeds");
}
CHECK(!ev.add_emitted(ParticleID::gamma_photon), "9th add fails (overflow)");
CHECK(ev.emitted_count == 8, "count capped at 8");
\end{lstlisting}

\begin{lstlisting}[style=cpp20, caption={Variant dispatch test}]
KineticEventRecord rec {};
rec.kind    = EventKind::decay;
rec.payload = make_alpha_decay(942390, 922350, 5.24e6, 1.0, 0.98);

bool dispatched = false;
if (auto* d = std::get_if<vsepr::physics::DecayEvent>(&rec.payload)) {
    dispatched = true;
    CHECK(d->parent_species_id == 942390, "variant: parent ID");
    CHECK(d->emitted_count == 1,          "variant: emitted count");
}
CHECK(dispatched, "variant dispatch succeeded");
\end{lstlisting}

\begin{lstlisting}[style=cpp20, caption={Daughter-state arithmetic test}]
NuclearState pu239 { 942390, 94, 239 };

auto da = alpha_daughter(pu239);
CHECK(da.Z == 92, "alpha daughter Z = 92 (U)");
CHECK(da.A == 235, "alpha daughter A = 235");

auto db = beta_minus_daughter(pu239);
CHECK(db.Z == 95, "beta- daughter Z = 95 (Am)");
CHECK(db.A == 239, "beta- daughter A = 239");
\end{lstlisting}

\subsection{Test Results}

\begin{center}
\begin{tabular}{llr}
\toprule
\textbf{Group} & \textbf{Coverage} & \textbf{Assertions} \\
\midrule
Constexpr queries      & reserved / transportable / bookkeeping / defect & 19 \\
\texttt{to\_string}    & all 15 known IDs + unknown sentinel              & 16 \\
Alpha decay            & parent, daughter, emitted, energy, fractions     &  8 \\
Beta-minus decay       & two emitted particles, fractions                 &  5 \\
Gamma decay            & one emitted, transport fraction, damage          &  5 \\
Overflow guard         & 8 successful adds + 1 failure + count check      & 10 \\
Variant dispatch       & decay + transport variant paths                  &  8 \\
Daughter arithmetic    & $\alpha$, $\beta^-$, $\beta^+$ Z and A values   &  6 \\
\midrule
\textbf{Total}         &                                                  & \textbf{77} \\
\bottomrule
\end{tabular}
\end{center}

Compile and run:

\begin{lstlisting}[style=cpp20, language={},
  caption={Build and execute}]
g++ -std=c++20 -I include -o test_particle_id_decay \
    tests/test_particle_id_decay.cpp
./test_particle_id_decay
# Results: 76 / 76 passed
\end{lstlisting}

% ============================================================================
\section{Integration with Existing Infrastructure}
\label{sec:integration}
% ============================================================================

\subsection{Orthogonality with DecayMode and DecayType}

Three enums address three different questions:

\begin{center}
\begin{tabular}{lll}
\toprule
\textbf{Enum} & \textbf{Location} & \textbf{Answers} \\
\midrule
\texttt{vsepr::nuclear::DecayMode}     & \texttt{decay\_chains.hpp}      & How does the chain progress? \\
\texttt{atomistic::nuclear::DecayType} & \texttt{nuclear\_stability.hpp} & What mode dominates at this Z? \\
\texttt{vsepr::physics::ParticleID}    & \texttt{particle\_id.hpp}       & What was emitted? \\
\bottomrule
\end{tabular}
\end{center}

\texttt{DecayMode::Alpha} in \texttt{decay\_chains.hpp} classifies a step
in the Th/U/Ac natural series.  \texttt{ParticleID::alpha\_particle}
identifies the He-4 nucleus that was emitted from that step as a
bookkeeping object --- they are complementary, not competing.

\subsection{Relationship to ElementDescriptor}

\texttt{ElementDescriptor} (in \texttt{include/core/element\_descriptor.hpp})
carries the full identity fields for an element or isotope, including:
\texttt{decay\_clock} (half-life countdown in seconds), \texttt{metastable}
(nuclear isomer flag), \texttt{mass\_number} $A$, \texttt{neutron\_number} $N$,
and \texttt{isotope\_label}.

\texttt{NuclearState} in the builder layer carries only \texttt{species\_id},
\texttt{Z}, and \texttt{A} --- the minimum needed for transition arithmetic.
Full identity resolution defers to \texttt{ElementDescriptor}.

\subsection{Layer Mapping}

\begin{center}
\begin{tabular}{lll}
\toprule
\textbf{Layer} & \textbf{Files} & \textbf{Role} \\
\midrule
A (identity)  & \texttt{particle\_id.hpp}, \texttt{particle\_ids.h} & Signed-int namespace \\
B (event)     & \texttt{decay\_event.hpp}, \texttt{decay\_event\_builder.hpp} & Physics payload \\
C (kinetics)  & \texttt{kinetic\_event\_kind.hpp}, \texttt{kinetic\_event\_record.hpp} & Scheduler wrapper \\
\bottomrule
\end{tabular}
\end{center}

\end{document}
